\documentclass[10pt,a4paper]{amsart}
\usepackage{amsmath,amssymb,amsthm,graphicx}
% Retained minimal preamble for the Stacks Project chapter varieties.tex.
% Verbatim excerpt of the chapter preamble: Xy-pic setup (native preamble lines 14-18), the native theorem
% declarations with the mathematical macros (native preamble lines 154-217), and the project own no-op
% definitions of the reference/slogan/history display environments (native wrapper.tex lines 1-3) that the
% counted statement uses for its slogan. Copied mechanically; nothing added, removed or redefined.
% For commutative diagrams we use Xy-pic
\usepackage[all]{xy}

% We use 2cell for 2-commutative diagrams.
\xyoption{2cell}
% Theorem environments.
%
\theoremstyle{plain}
\newtheorem{theorem}[subsection]{Theorem}
\newtheorem{proposition}[subsection]{Proposition}
\newtheorem{lemma}[subsection]{Lemma}

\theoremstyle{definition}
\newtheorem{definition}[subsection]{Definition}
\newtheorem{example}[subsection]{Example}
\newtheorem{exercise}[subsection]{Exercise}
\newtheorem{situation}[subsection]{Situation}

\theoremstyle{remark}
\newtheorem{remark}[subsection]{Remark}
\newtheorem{remarks}[subsection]{Remarks}

\numberwithin{equation}{subsection}

% Macros
%
\def\lim{\mathop{\mathrm{lim}}\nolimits}
\def\colim{\mathop{\mathrm{colim}}\nolimits}
\def\Spec{\mathop{\mathrm{Spec}}}
\def\Hom{\mathop{\mathrm{Hom}}\nolimits}
\def\Ext{\mathop{\mathrm{Ext}}\nolimits}
\def\SheafHom{\mathop{\mathcal{H}\!\mathit{om}}\nolimits}
\def\SheafExt{\mathop{\mathcal{E}\!\mathit{xt}}\nolimits}
\def\Sch{\mathit{Sch}}
\def\Mor{\mathop{\mathrm{Mor}}\nolimits}
\def\Ob{\mathop{\mathrm{Ob}}\nolimits}
\def\Sh{\mathop{\mathit{Sh}}\nolimits}
\def\NL{\mathop{N\!L}\nolimits}
\def\CH{\mathop{\mathrm{CH}}\nolimits}
\def\proetale{{pro\text{-}\acute{e}tale}}
\def\etale{{\acute{e}tale}}
\def\QCoh{\mathit{QCoh}}
\def\Ker{\mathop{\mathrm{Ker}}}
\def\Im{\mathop{\mathrm{Im}}}
\def\Coker{\mathop{\mathrm{Coker}}}
\def\Coim{\mathop{\mathrm{Coim}}}

% Boxtimes
%
\DeclareMathSymbol{\boxtimes}{\mathbin}{AMSa}{"02}

%
% Macros for moduli stacks/spaces
%
\def\QCohstack{\mathcal{QC}\!\mathit{oh}}
\def\Cohstack{\mathcal{C}\!\mathit{oh}}
\def\Spacesstack{\mathcal{S}\!\mathit{paces}}
\def\Quotfunctor{\mathrm{Quot}}
\def\Hilbfunctor{\mathrm{Hilb}}
\def\Curvesstack{\mathcal{C}\!\mathit{urves}}
\def\Polarizedstack{\mathcal{P}\!\mathit{olarized}}
\def\Complexesstack{\mathcal{C}\!\mathit{omplexes}}
% \Pic is the operator that assigns to X its picard group, usage \Pic(X)
% \Picardstack_{X/B} denotes the Picard stack of X over B
% \Picardfunctor_{X/B} denotes the Picard functor of X over B
\def\Pic{\mathop{\mathrm{Pic}}\nolimits}
\def\Picardstack{\mathcal{P}\!\mathit{ic}}
\def\Picardfunctor{\mathrm{Pic}}
\def\Deformationcategory{\mathcal{D}\!\mathit{ef}}
\newenvironment{reference}{}{}
\newenvironment{slogan}{}{}
\newenvironment{history}{}{}
\usepackage{cleveref}
\crefname{originallemma}{original lemma}{original lemmas}
\Crefname{originallemma}{Original lemma}{Original lemmas}
\makeatletter
\newlabel{lemma-kunneth}{{0BED}{source}{}{sourcepointer-1}{}}
\newlabel{lemma-kunneth@cref}{{[originallemma][0][]0BED}{[1][0][]source}{}{}{}}
\makeatother
\crefname{originalreference}{original reference}{original references}
\Crefname{originalreference}{Original reference}{Original references}
\makeatletter
\newlabel{coherent-lemma-hypercoverings}{{01XN}{source}{}{sourcepointer-2}{}}
\newlabel{coherent-lemma-hypercoverings@cref}{{[originalreference][0][]01XN}{[1][0][]source}{}{}{}}
\makeatother
\crefname{originalreference}{original reference}{original references}
\Crefname{originalreference}{Original reference}{Original references}
\makeatletter
\newlabel{cohomology-lemma-before-Leray}{{01EZ}{source}{}{sourcepointer-3}{}}
\newlabel{cohomology-lemma-before-Leray@cref}{{[originalreference][0][]01EZ}{[1][0][]source}{}{}{}}
\makeatother
\crefname{originalreference}{original reference}{original references}
\Crefname{originalreference}{Original reference}{Original references}
\makeatletter
\newlabel{coherent-lemma-vanishing-nr-affines-quasi-separated}{{071L}{source}{}{sourcepointer-4}{}}
\newlabel{coherent-lemma-vanishing-nr-affines-quasi-separated@cref}{{[originalreference][0][]071L}{[1][0][]source}{}{}{}}
\makeatother
\crefname{originalreference}{original reference}{original references}
\Crefname{originalreference}{Original reference}{Original references}
\makeatletter
\newlabel{cohomology-lemma-quasi-separated-cohomology-colimit}{{01FF}{source}{}{sourcepointer-5}{}}
\newlabel{cohomology-lemma-quasi-separated-cohomology-colimit@cref}{{[originalreference][0][]01FF}{[1][0][]source}{}{}{}}
\makeatother
\title{The Kunneth formula for quasi-compact quasi-separated schemes via square-zero thickenings}
\author{The Stacks Project Authors}
\date{stacks-project:varieties.tex:local-snapshot-20261011; GNU FDL 1.2 or later}
\begin{document}
\maketitle
\noindent\textit{Editorial note.} All delivered mathematical text is the author's own, copied verbatim from the pinned Stacks Project chapter file varieties.tex; no statement or proof step was rewritten, reordered or removed.\par\medskip
\noindent\textit{Editorial note.} The counted span cites five cross-chapter labels, each recorded in assembly.external\_references with its canonical tag as a navigation pointer.\par\medskip
\noindent\textit{Editorial note.} The proof cites no bibliography item inside the counted span.\par\medskip
\noindent\textit{Allowed standard prerequisites (editorial).}\begin{itemize}
\item Cohomology of a square-zero thickening and the associated exact sequences. Stacks Project, Cohomology of Schemes (cohomology of infinitesimal thickenings). Provides the exact sequences on which the comparison of the thickenings is based.
\item The affine-diagonal Kunneth formula. Stacks Project, Varieties, Lemma lemma-kunneth (tag 0BKC). Handles the separated/affine-diagonal case, which the author directs the reader to for the sharper statement about the maps.
\end{itemize}\medskip
\ifdefined\hypertarget\hypertarget{sourcepointer-1}{}\fi
\noindent\textit{Source reference.} Original lemma 0BED: sources/stacks-chow/tags.txt, source lines 14346--14346. Cross-chapter standard prerequisite cited by the author inside the counted span; it is used as a pointer and its own proof is not part of this item.\par\medskip
\ifdefined\hypertarget\hypertarget{sourcepointer-2}{}\fi
\noindent\textit{Source reference.} Original reference 01XN: sources/stacks-chow/tags.txt, source lines 2526--2526. Cross-chapter standard prerequisite cited by the author inside the counted span; it is used as a pointer and its own proof is not part of this item.\par\medskip
\ifdefined\hypertarget\hypertarget{sourcepointer-3}{}\fi
\noindent\textit{Source reference.} Original reference 01EZ: sources/stacks-chow/tags.txt, source lines 1852--1852. Cross-chapter standard prerequisite cited by the author inside the counted span; it is used as a pointer and its own proof is not part of this item.\par\medskip
\ifdefined\hypertarget\hypertarget{sourcepointer-4}{}\fi
\noindent\textit{Source reference.} Original reference 071L: sources/stacks-chow/tags.txt, source lines 8995--8995. Cross-chapter standard prerequisite cited by the author inside the counted span; it is used as a pointer and its own proof is not part of this item.\par\medskip
\ifdefined\hypertarget\hypertarget{sourcepointer-5}{}\fi
\noindent\textit{Source reference.} Original reference 01FF: sources/stacks-chow/tags.txt, source lines 1868--1868. Cross-chapter standard prerequisite cited by the author inside the counted span; it is used as a pointer and its own proof is not part of this item.\par\medskip
\begin{lemma}
\label{lemma-kunneth-general}
Let $k$ be a field. Let $X$ and $Y$ be schemes over $k$ and
let $\mathcal{F}$, resp.\ $\mathcal{G}$ be a quasi-coherent
$\mathcal{O}_X$-module, resp.\ $\mathcal{O}_Y$-module.
Then we have a canonical isomorphism
$$
H^n(X \times_{\Spec(k)} Y, \text{pr}_1^*\mathcal{F}
\otimes_{\mathcal{O}_{X \times_{\Spec(k)} Y}} \text{pr}_2^*\mathcal{G}) =
\bigoplus\nolimits_{p + q = n}
H^p(X, \mathcal{F}) \otimes_k H^q(Y, \mathcal{G})
$$
provided $X$ and $Y$ are quasi-compact and quasi-separated.
\end{lemma}

\begin{proof}
If $X$ and $Y$ are separated or more generally have affine diagonal, then
please see Lemma \ref{lemma-kunneth} for ``better'' proof
(the feature it has over this proof is that it identifies the
maps as pullbacks followed by cup products).
Let $X'$, resp.\ $Y'$ be the infinitesimal thickening of $X$, resp.\ $Y$
whose structure sheaf is
$\mathcal{O}_{X'} = \mathcal{O}_X \oplus \mathcal{F}$,
resp.\ $\mathcal{O}_{Y'} = \mathcal{O}_Y \oplus \mathcal{G}$ where
$\mathcal{F}$, resp.\ $\mathcal{G}$ is an ideal of square zero.
Then
$$
\mathcal{O}_{X' \times Y'} =
\mathcal{O}_{X \times Y}
\oplus \text{pr}_1^*\mathcal{F}
\oplus \text{pr}_2^*\mathcal{G}
\oplus \text{pr}_1^*\mathcal{F}
\otimes_{\mathcal{O}_{X \times Y}} \text{pr}_2^*\mathcal{G}
$$
as sheaves on $X \times Y$. In this way we see that it suffices to
prove that
$$
H^n(X \times Y, \mathcal{O}_{X \times Y}) =
\bigoplus\nolimits_{p + q = n}
H^p(X, \mathcal{O}_X) \otimes_k H^q(Y, \mathcal{O}_Y)
$$
for any pair of quasi-compact and quasi-separated schemes over $k$.
Some details omitted.

\medskip\noindent
To prove this statement we use cohomology and base change in the form
of Cohomology of Schemes, Lemma \ref{coherent-lemma-hypercoverings}.
This lemma tells us there exists a bounded below complex of $k$-vector spaces,
i.e., a complex $\mathcal{K}^\bullet$ of quasi-coherent modules on $\Spec(k)$,
which universally computes the cohomology of $Y$ over $\Spec(k)$.
In particular, we see that
$$
R\text{pr}_{1, *}(\mathcal{O}_{X \times Y}) \cong
(X \to \Spec(k))^*\mathcal{K}^\bullet
$$
in $D(\mathcal{O}_X)$. Up to homotopy the complex $\mathcal{K}^\bullet$
is isomorphic to $\bigoplus_{q \geq 0} H^q(Y, \mathcal{O}_Y)[-q]$ because
this is true for every complex of vector spaces over a field.
We conclude that
$$
R\text{pr}_{1, *}(\mathcal{O}_{X \times Y}) \cong
\bigoplus\nolimits_{q \geq 0}
H^q(Y, \mathcal{O}_Y)[-q] \otimes_k \mathcal{O}_X
$$
in $D(\mathcal{O}_X)$. Then we have
\begin{align*}
R\Gamma(X \times Y, \mathcal{O}_{X \times Y})
& =
R\Gamma(X, R\text{pr}_{1, *}(\mathcal{O}_{X \times Y})) \\
& =
R\Gamma(X, \bigoplus\nolimits_{q \geq 0}
H^q(Y, \mathcal{O}_Y)[-q] \otimes_k \mathcal{O}_X) \\
& =
\bigoplus\nolimits_{q \geq 0}
R\Gamma(X,  H^q(Y, \mathcal{O}_Y) \otimes \mathcal{O}_X)[-q] \\
& =
\bigoplus\nolimits_{q \geq 0}
R\Gamma(X, \mathcal{O}_X) \otimes_k H^q(Y, \mathcal{O}_Y)[-q] \\
& =
\bigoplus\nolimits_{p, q \geq 0}
H^p(X, \mathcal{O}_X)[-p] \otimes_k H^q(Y, \mathcal{O}_Y)[-q]
\end{align*}
as desired. The first equality by Leray for $\text{pr}_1$
(Cohomology, Lemma \ref{cohomology-lemma-before-Leray}).
The second by our decomposition of the total direct image given above.
The third because cohomology always commutes with finite direct sums
(and cohomology of $Y$ vanishes in sufficiently large degree by
Cohomology of Schemes, Lemma
\ref{coherent-lemma-vanishing-nr-affines-quasi-separated}).
The fourth because cohomology on $X$ commutes with infinite
direct sums by Cohomology, Lemma
\ref{cohomology-lemma-quasi-separated-cohomology-colimit}.
The final equality by our remark on the derived category
of a field above.
\end{proof}

\end{document}
